\section{Scalar three-\texorpdfstring{$Z$}{Z} input for the lagwise route}
\label{sec:scalar-hlz-core}

This section records only the scalar Heap--Li--Zhao input used in the attempted
lagwise route.  The arithmetic lift from the physical one-$\chi$ carrier to
this untwisted scalar diagonal is a separate issue; see
\texttt{doc/lagwise\_hlz\_gap\_audit.md}.  The packet variable has not yet been
converted into a one-sided prefix coordinate.  The only packet information
admitted in the scalar lemma below is a bounded-support Mellin weight.

For the invariant outer specialization put
\begin{equation}
\boxed{
Z_{\alpha,\beta,\gamma,1,1}
:=\frac{\zeta(1+\alpha+\beta)\zeta(1+\beta+\gamma)}
        {\zeta(1+\beta)}.
}
\label{eq:scalar-HLZ-Euler-object}
\end{equation}
This is the $h=k=1$ Euler object in the scalar HLZ diagonal.  In particular,
there is no auxiliary outer arithmetic coordinate in this section.

Fix constants
\[
0<r_1<r_2<r_3,\qquad r_{j+1}-r_j\ge r_*>0,
\]
and put
\[
\mathcal T_{j,L}:=\{z:|z|=r_j/L\},\qquad
\mathcal T_L:=\mathcal T_{1,L}\times\mathcal T_{2,L}\times\mathcal T_{3,L}.
\]
Thus
\[
|z_i-z_j|\ge r_*/L\qquad(i\ne j)
\]
on the separated Cauchy torus.  For $\mathbf z=(z_1,z_2,z_3)\in\mathcal T_L$
define
\begin{equation}
\boxed{
H_{\mathbf z}(u):=
\exp(u^2)
\frac{(2u-z_2+z_1)(2u-z_2+z_3)}
     {(z_1-z_2)(z_3-z_2)}.
}
\label{eq:scalar-HLZ-auxiliary}
\end{equation}
Then
\begin{equation}
H_{\mathbf z}(0)=1,
\qquad
H_{\mathbf z}\!\left(\frac{z_2-z_1}{2}\right)=0,
\qquad
H_{\mathbf z}\!\left(\frac{z_2-z_3}{2}\right)=0.
\label{eq:scalar-HLZ-auxiliary-zeros}
\end{equation}
The separated torus is fixed throughout; no analytic continuation of
$H_{\mathbf z}$ through the collision locus is used.

\begin{proposition}[Exact three-face Cauchy residue identity]
\label{prop:scalar-HLZ-three-face}
Let
\[
\mathfrak Z_{\alpha,\beta,\gamma}(X)
:=Z_{\alpha,\beta,\gamma,1,1}
+X^{-\alpha-\beta}Z_{-\beta,-\alpha,\gamma,1,1}
+X^{-\beta-\gamma}Z_{\alpha,-\gamma,-\beta,1,1}.
\]
With $\Delta(\mathbf z)=\prod_{i<j}(z_j-z_i)$, one has the exact algebraic
identity
\begin{equation}
\boxed{
\begin{aligned}
\mathfrak Z_{\alpha,\beta,\gamma}(X)
={}&-\frac12X^{-(\alpha+\beta+\gamma)/2}
\frac1{(2\pi i)^3}
\iiint_{\mathcal T_L}
Z_{z_1,-z_2,z_3,1,1}\\
&\qquad\times
\frac{\Delta(\mathbf z)^2X^{(z_1-z_2+z_3)/2}}
{\prod_{j=1}^3(z_j-\alpha)(z_j+\beta)(z_j-\gamma)}
\,d\mathbf z .
\end{aligned}}
\label{eq:scalar-HLZ-three-face}
\end{equation}
The identity is a finite three-variable Cauchy residue calculation.  It is
consistent with the three principal faces occurring in the completed HLZ
expansion, but it is proved here directly rather than attributed to a formula
number in Heap--Li--Zhao.
\end{proposition}

\begin{proof}
Put
\[
P(z):=(z-\alpha)(z+\beta)(z-\gamma),
\]
whose roots are
\[
r_1=\alpha,\qquad r_2=-\beta,\qquad r_3=\gamma.
\]
The Vandermonde square in \eqref{eq:scalar-HLZ-three-face} vanishes whenever two
of the three Cauchy variables take the same root.  Hence only the six
permutations of $(r_1,r_2,r_3)$ contribute.  For every such permutation,
\[
\frac{\Delta(r_1,r_2,r_3)^2}
     {P'(r_1)P'(r_2)P'(r_3)}=-1,
\]
because for a monic cubic
$P'(r_i)=\prod_{j\ne i}(r_i-r_j)$ and the product of the three derivatives is
$-\Delta(r_1,r_2,r_3)^2$.

The remaining Euler factor and the height monomial are symmetric under
$z_1\leftrightarrow z_3$.  Thus the six residues pair into three equal pairs,
and the prefactor $-1/2$ turns each pair into one copy of the corresponding
face.  If the middle variable is $z_2=-\beta$, the outer height exponent
cancels and the first face is $Z_{\alpha,\beta,\gamma,1,1}$.  If
$z_2=\alpha$, the total height exponent is $-\alpha-\beta$, giving
$X^{-\alpha-\beta}Z_{-\beta,-\alpha,\gamma,1,1}$.  If
$z_2=\gamma$, it is $-\beta-\gamma$, giving
$X^{-\beta-\gamma}Z_{\alpha,-\gamma,-\beta,1,1}$.  Their sum is exactly
\eqref{eq:scalar-HLZ-three-face}.
\end{proof}

\begin{lemma}[Uniform scalar HLZ diagonal for bounded-support Mellin weights]
\label{lem:scalar-lagweight-HLZ}
Put $\log_3T:=\log\log\log T$ and
\[
D_{\mathbf z}(u):=
\frac{\zeta(1+z_1-z_2+2u)\zeta(1-z_2+z_3+2u)}
     {\zeta(1-z_2+2u)}.
\]
Let $c_+>0$ be fixed so that this Euler product is absolutely convergent on
$\Re u=c_+/L$.  Suppose
\begin{equation}
W_{\rho,p,\mathbf z,T}(u)
=\int_0^\rho a_{\rho,p,\mathbf z,T}(t)e^{Lut}\,dt
\label{eq:scalar-HLZ-admissible-weight}
\end{equation}
(or a fixed finite product of such factors), where $0\le\rho\le C_*$ and, for
every fixed derivative order used below,
\begin{equation}
\left\|\partial_t^m\partial_\rho^r\partial_p^j
 a_{\rho,p,\mathbf z,T}\right\|_{L^1(0,C_*)}
\ll L^{C_{m,r,j}}
\label{eq:scalar-HLZ-weight-bounds}
\end{equation}
uniformly on $\mathcal T_L$.  Then, pointwise on the separated torus,
\begin{equation}
\boxed{
W(0)Z_{\mathbf z,1,1}
=\frac1{2\pi i}\int_{c_+/L-i\infty}^{c_+/L+i\infty}
\frac{H_{\mathbf z}(u)W(u)}u T^{2u}D_{\mathbf z}(u)\,du
+\mathcal E_{\rm diag}(\rho,p,\mathbf z),
}
\label{eq:scalar-HLZ-diagonal}
\end{equation}
where, for every prescribed $A>0$ and every fixed $r,j$ required later,
\begin{equation}
\boxed{
\sup_{\mathbf z\in\mathcal T_L}
|\partial_\rho^r\partial_p^j\mathcal E_{\rm diag}|
\ll_A L^{-A+C_{r,j}}.
}
\label{eq:scalar-HLZ-error}
\end{equation}
The same estimate survives the finite Cauchy residues in
Proposition~\ref{prop:scalar-HLZ-three-face} after increasing the requested
saving exponent.
\end{lemma}

\begin{proof}
This is the scalar diagonal contour mechanism of Heap--Li--Zhao~\cite{HLZ},
specialized to $h=k=1$ and multiplied by the admissible analytic weight $W$.
On the strip
\[
-1/\log_3T\le\Re u\le c_+/L
\]
the bounded support in \eqref{eq:scalar-HLZ-admissible-weight} and
\eqref{eq:scalar-HLZ-weight-bounds} give only fixed powers of $L$.  The factor
$H_{\mathbf z}(u)$ contributes a fixed polynomial in $L(1+|u|)$ times the
Gaussian $e^{-\Im(u)^2}$.  Its two zeros in
\eqref{eq:scalar-HLZ-auxiliary-zeros} cancel the shifted zeta poles
$2u=z_2-z_1$ and $2u=z_2-z_3$.

Truncate at $|\Im u|=B\sqrt{\log L}$ and move the line to
$\Re u=-1/\log_3T$, as in the scalar HLZ diagonal argument.  The left edge is
\[
\ll L^C T^{-2/\log_3T}\log_3T,
\]
while the two horizontal edges and Gaussian tails are $O_A(L^{-A})$.  Since
\[
T^{-2/\log_3T}L^D\log_3T\ll_{A,D}L^{-A}
\]
for every fixed $A,D$, the complete error has arbitrary logarithmic saving.
Differentiating in $\rho$ produces only endpoint jets and fixed derivatives of
the amplitude; differentiating in $p$ is covered directly by
\eqref{eq:scalar-HLZ-weight-bounds}.  Thus the same contour proof yields
\eqref{eq:scalar-HLZ-error}.  Finally the separated Cauchy denominators cost
only fixed powers of $L$, which are absorbed by requesting a larger scalar
saving exponent before the finite residues are taken.
\end{proof}

\begin{lemma}[Local-height rescaling of the scalar diagonal]
\label{lem:scalar-HLZ-local-height}
Let $X(\tau)=\tau/(2\pi)$ with $\tau\asymp T$.  In
\eqref{eq:scalar-HLZ-diagonal}, the replacement
\[
T^{2u}\longmapsto X(\tau)^{2u}
\]
does not alter the residue at $u=0$ and preserves the error class
\eqref{eq:scalar-HLZ-error}, including the first two fixed endpoint and
curvature derivatives.
\end{lemma}

\begin{proof}
Multiply the auxiliary weight by
\[
R_\tau(u):=\left(\frac{X(\tau)}T\right)^{2u}.
\]
Since $X(\tau)/T$ remains in a fixed compact subset of $(0,\infty)$,
$R_\tau$ and every fixed height derivative have polynomial growth in $u$ on the
same strip, hence preserve the Gaussian contour estimates.  Since
$R_\tau(0)=1$, the $u=0$ residue is unchanged.  The same argument after the
fixed endpoint and curvature derivatives proves the differentiated statement.
\end{proof}

\begin{remark}[Scope and current obstruction]
The scalar statements above concern the untwisted Euler diagonal only.  The
physical packet problem still contains the additive stationary carrier
$e(-e^d\ell)$ after fixing the lag.  No assertion in this section removes that
factor.  Consequently Proposition~\ref{prop:lagwise-packet-HLZ} requires an
additional arithmetic bridge which is presently open; see
\texttt{doc/lagwise\_hlz\_gap\_audit.md}.  In particular, the scalar lemma alone
does not provide the theorem-facing packet lift.
\end{remark}
